\section{Технологийн судалгаа}

\section{Сэдвийн судлагдсан байдал}
\label{sec:literature-review}
Энэ хэсэгт микрофронтэнд архитектур зохиомж, түүний янз бүрийн шинж чанар, хэрэгжүү- лэлт болон өмнө нь хийгдсэн судалгааны ажлуудын тоймыг хүргэх болно. Одоо байгаа ном зохиол, судалгааг тоймлохын гол зорилго нь микрофронтэнд архитектур зохиомжтой холбоотой орчин үеийн тэргүүлэх шийдлүүдийг тодорхойлох,  холбогдох хэрэгсэл, багц, санг судлахад оршино.
\subsection{Судалгааны ажлын хайлт}
\hspace*{1cm}Холбогдох өгүүллүүдийг олохын тулд Scopus\footnote{Scopus нь шинжлэх ухааны сэтгүүл, ном, хурлын бүтээлийн хураангуй болон эшлэлийн мэдээллийг нэгтгэсэн олон улсын хамгийн том өгөгдлийн сан юм.}, DiVA\footnote{DiVA нь магистр болон докторын зэрэг хамгаалах ажил, судалгааны бүтээл, нийтлэл, сэтгүүл зэрэг шинжлэх ухааны бүтээлүүдийн цуглуулга бүхий сан юм.} өгөгдлийн санг ашигласан. Хайлтын утгад микрофронтэнд архитектур зохиомжийн нэршлийн өөр өөр хувилбаруудыг ашигласан болно. Үүнд: \verb|micro-frontend| \verb|OR| \verb|microfrontend| \verb|OR| \verb|"micro frontend*"**| гэсэн хайлтын утгаар Scopus өгөгдлийн сангаас нийт 17 илэрц, \verb|microfrontend| \verb|OR| \verb|micro frontend| \verb|OR| \verb|micro-frontend| гэсэн хайлтын утгаар DiVA сангаас 2 илэрц олдсон. Өгүүллийн гарчиг, хураангуй, эсвэл түлхүүр үгсэд хайлтын утга орсон бүтээлүүдийг сонгож авсан. Олдсон өгүүллүүдийг уншиж танилцсаны дараа судалгааны арга зүй, сэдвийн хамаарал, сэдвийг хэрхэн хөндсөн байдал зэрэгт нь үндэслэн 5 өгүүллийг тоймлон үзэхээр сонгосон. Эдгээрийн 1 нь DiVA сангаас, үлдсэн 4 нь Scopus өгөгдлийн сангаас авагдсан болно.

\subsection{Сонгож авсан судалгааны ажил}
\hspace*{1cm}[*] дугаартай бүтээлд зохиогчид нь одоо ашиглагдаж буй QL4POMR бэкэнд программ хангамжийн шийдэл дээр тулгуурлан эмч, өвчтөн хоорондын харилцан яриаг загварчлах зориулалттай чатботыг хэрхэн хэрэгжүүлсэн талаар тайлбарласан байдаг. Энэхүү харилцан ярианы систем нь patientMFE болон physicianMFE гэж нэрлэгдсэн хоёр микрофронтэнд аппли- кэйшнээс бүрдэнэ. Уг бүтээлд микрофронтэнд хэрэгжүүлэхэд ашиглагддаг Piral, SystemJS, Single SPA, Webpack Module Federation, Bit.Dev зэрэг зарим хэрэгслүүдийг дурджээ. Хэрэг- жүүлэлтийн техникийн шийдэл нь зөвхөн тухайн системийн тохиргоонд зориулан ашигласан технологиудаар хязгаарлагдсан байна.
\newline
\hspace*{1cm}[*] дугаартай бүтээлд Стефановска болон Трайковик нар домэйнд суурилсан зохиомж нь микрофронтэнд архитектурын үндсэн гол ойлголт болохын хувьд түүний тоймоор судал- гаагаа эхлүүлсэн байна. Тэд одоо байгаа техникийн шийдлүүдийг чанарын үзүүлэлтүүд ашиглан үнэлжээ. Үнэлгээ хийсний дараа тэд кодын дахин ашиглалтыг чухалчилсан өөрсдийн техникийн шийдлээ санал болгосон. Энэхүү техникийн шийдэл нь 2 ажлын явцаас бүрдсэн. Эхний ажлын явцад вэб компонент нь фронтэнд аппликэйшнийг бүхэлд нь агуулах бүрхүүл байдлаар хэрхэн ажилладгийг тайлбарласан. Хоёр дахь ажлын явц нь модулийн шинж чанарт төвлөрсөн бөгөөд үүнийг Module Federation Plugin ашиглан хэрэгжүүлсэн байна. Тэдний дурдсан Module Federation Plugin-ийг ашигласнаар олж авах боломжтой нэг шинж чанар нь өргөжүүлэх боломж юм. Учир нь микрофронтэнд аппликэйшнүүдийг статик JavaScript файл хэлбэрээр татаж авдаг байна. Санал болгосон хэрэгжүүлэлтийн үр дүнд аппликэйшний хөгжүүлэлтээс үйлдвэрлэл рүү шилжих амьдралын мөчлөгийн хугацаа илүү хурдан болсон байна.
\newline
\hspace*{1cm}[*] дугаартай бүтээлд вэб хөгжүүлэлт дэх микрофронтэнд архитектурын ойлголтуудыг нарийвчлан судалсан байдаг. Зохиогчид үүнийг микросервис архитектурын дараагийн салш- гүй, зүй ёсны хөгжил гэж үзсэн. Уг судалгаанд вэб болон бэкэнд хөгжүүлэлтийн уламжлалт болон микросервис архитектурын хандлагуудыг харьцуулж, микрофронтэндийн өнөөгийн туршлагуудыг судалжээ. Хөгжүүлэгчдийн туршлагад үндэслэн кэйс судалгаа хийсэн бөгөөд чиглүүлэлт, статик файл дамжуулалт, репозиторын зохион байгуулалт зэрэгт тулгардаг сорил- туудыг онцолсон байна. Микрофронтэнд нь том баг, нарийн төвөгтэй, томоохон хэмжээний төслүүдэд илүү тохиромжтой гэж дүгнэсэн. Харин цөөн хөгжүүлэгчтэй багийн хувьд уг архитектурыг дэмжиж ажиллахад нэмэлт хүчин чармайлт шаарддаг бөгөөд төслийн бүтэц, хөгжүүлэлтийн үйл явцыг нэлээд төвөгтэй болгодог байна. Мөн уг архитектурын давуу талыг бүрэн ашиглахын тулд хөгжүүлэгчдээс тодорхой хэмжээний туршлага шаарддаг төдий- гүй нийгэмлэгийн зүгээс үүнийг бүрэн дэмжих "бэлэн шийдлүүд" хомс байгааг дурджээ.
\newline
\hspace*{1cm}[*] дугаартай өгүүллийн зорилго нь хөгжүүлэлтийн багийн бүтээмжийг нэмэгдүүлэх, нөөцийн зарцуулалтыг сайжруулахад чиглэсэн платформын шийдлийг танилцуулахад оршино. Уг платформ нь нүсэр, нийлмэл программ хангамжийн хөгжүүлэлтийг төвлөрсөн бус арга зүйгээр хөнгөвчлөх бөгөөд микросервис болон микрофронтэндийг үүсгэх, ашиглах үйл явцыг дэмждэг. Энэхүү судалгаа нь микросервис болон микрофронтэнд архитектурын давуу тал болон сорилтуудыг ерөнхийд нь тоймлон харуулсан. Зэрэгцээ хөгжүүлэлт болон өргөтгөх боломж нь эдгээр архитектурын хамгийн том давуу тал бөгөөд ингэснээр нөөцийг илүү үр дүнтэй ашиглах боломж бүрддэг байна. Мөн фронтэндийн хэсгүүдийг алхам алхмаар шинэч- лэх боломж нь уламжлалт цул архитектуртай харьцуулахад хөгжүүлэгчийн бүтээмжийг эрс нэмэгдүүлдэг. Эцэст нь, зохиогчид бүтээгч, хэрэглэгч болон администраторуудад зориулсан гурван хэсэгтэй платформыг санал болгосон нь системийг хөгжүүлэхэд тулгардаг хүндрэл- үүдийг илүү сайн ойлгох, цаашлаад платформын үйл ажиллагааг сайжруулахад чухал түлхэц болжээ.
\newline
\hspace*{1cm}[*] дугаартай бүтээлд микрофронтэндийн талаарх хайгуулын судалгаа хийж, хөгжүүлэгч- дээс ярилцлага авсан байна. Энэхүү судалгаа нь системийн өөрчлөгдөх боломж болон ерөнхий- лөн ашиглагдах байдал зэрэг чанарын үзүүлэлтүүд дээр төвлөрч, микрофронтэндийн давуу тал, тулгарч буй асуудлуудыг тодруулахыг зорьжээ. Судалгааны хүрээнд нэг нь SPA, нөгөө нь микрофронтэнд архитектур ашигласан хоёр өөр фронтэнд аппликэйшн хөгжүүлж туршсан байна. Үүний дараа уламжлалт болон микрофронтэнд архитектурын аль алинд нь туршлагатай хөгжүүлэгчдээс ярилцлага авч, өөрчлөгдөх боломжийг үнэлэхийн тулд хэд хэдэн хэмжүүр ашиглажээ. Эдгээр өгөгдөл дээр үндэслэн микрофронтэнд архитектурыг сонгох үед анхаарах ёстой эрсдэлүүдийг хэлэлцсэн байна. Эцсийн дүгнэлтдээ Монтелиус ярилцлагын үр дүнг туйлын үнэн гэхээсээ илүүтэйгээр өөр өнцгөөс харсан хандлага гэж үзэх нь зүйтэйг онцолсон бөгөөд турших загваруудаас гарсан үр дүнг бүх SPA болон микрофронтэндүүдэд шууд хамаа- руулан ерөнхийлөн дүгнэх боломжгүйг сануулжээ.

\subsection{Дэд бүлгийн дүгнэлт}

Дээрх судлагдсан бүтээлүүдийн ихэнх нь программ хангамжийн чанарыг ISO 25010 стан- дартын дагуу үнэлсэн байсан нь энэхүү дипломын ажлын арга зүйтэй ижил төстэй байна. Зарим судалгаанууд микрофронтэнд архитектурыг онолын түвшинд тайлбарлаж, уламжлалт цул архитектураас шилжих загвар санал болгожээ. Түүнчлэн SPA болон микрофронтэнд архитек- турыг хооронд нь харьцуулах замаар давуу болон сул талуудыг гаргаж ирсэн байна. Микрофронтэнд архитектурын шийдлийг санал болгож буй судалгаануудын ихэнх нь Module Federation Plugin-ийг ашигласан байв.